\addcontentsline{toc}{section}{СПИСОК ИСПОЛЬЗОВАННЫХ ИСТОЧНИКОВ}

\begin{thebibliography}{9}
	
	\bibitem{python} Лутц, М. Изучаем Python / М. Лутц. – Санкт-Петербург: Питер, 2020. – 1600 с. – ISBN 978-5-4461-1452-0. – Текст: непосредственный.
	
	\bibitem{tkinter} Либерман, М. Tkinter. Создание графических интерфейсов на Python / М. Либерман. – Москва: Вильямс, 2019. – 384 с. – ISBN 978-5-8459-2391-2. – Текст: непосредственный.
	
	\bibitem{customtkinter} Рейнольдс, А. Modern Python GUI with CustomTkinter / A. Reynolds. – 2022. – URL: https://github.com/TomSchimansky/CustomTkinter
	
	\bibitem{markov} Марков, А. О произведениях в математике / А. Марков. – Москва: Наука, 1970. – 320 с. – Текст: непосредственный.
	
	\bibitem{prod_systems} Виноградов, В. Продукционные системы и их применение в вычислительной технике / В. Виноградов. – Санкт-Петербург: Питер, 2015. – 256 с. – ISBN 978-5-496-01502-3. – Текст: непосредственный.
	
	\bibitem{algorithms} Кормен, Т., Лейзерсон, Ч., Ривест, Р., Штайн, К. Алгоритмы: построение и анализ / Т. Кормен и др. – Москва: Вильямс, 2013. – 1292 с. – ISBN 978-5-8459-1806-2. – Текст: непосредственный.
	
	\bibitem{logging} Миллер, М. Логирование в Python: практическое руководство / М. Миллер. – Москва: ДМК Пресс, 2020. – 200 с. – ISBN 978-5-97060-626-4. – Текст: непосредственный.
	
	\bibitem{gui_design} Шнайдер, Э. Разработка графических интерфейсов: принципы и практика / Э. Шнайдер. – Санкт-Петербург: Питер, 2018. – 336 с. – ISBN 978-5-496-01734-8. – Текст: непосредственный.
	
	\bibitem{statistics} Харрис, М. Статистика для программистов на Python / М. Харрис. – Москва: Питер, 2021. – 312 с. – ISBN 978-5-4461-2410-0. – Текст: непосредственный.
	
	\bibitem{python_official} Python Software Foundation. Python 3 Documentation / PSF. – URL: https://docs.python.org/3/
	
\end{thebibliography}
